\documentclass[11pt]{article}
\usepackage{xltxtra}
\usepackage{bookmark}
\usepackage{hyperref}
\hypersetup{hidelinks}
\usepackage{url}
\urlstyle{tt}
\usepackage{multicol}
\usepackage{xcolor}
\usepackage{calc}
\usepackage{graphicx}
\usepackage{tikz}
\usetikzlibrary{calc}
\usepackage{fontspec}
\usepackage{xeCJK}
\usepackage{relsize}
\usepackage{xspace}
\usepackage{fontawesome5}
\usepackage{titlesec}
\usepackage{enumitem}
\usepackage{siunitx}
\usepackage{amssymb}
\usepackage{tabularx}
\usepackage{multicol}
\usepackage{fontspec}
\usepackage{fancybox}
\usepackage{float}

% 一些小设置
\CJKsetecglue{}
\protected\def\Cpp{{C\nolinebreak[4]\hspace{-.05em}\raisebox{.28ex}{\relsize{-1}++}}\xspace}
\setlength{\parindent}{0pt}
\pagenumbering{gobble}
\setlist[itemize]{topsep=0em, leftmargin=*}
\setlist[enumerate]{topsep=0em, leftmargin=*}

% --- 华中科技大学配色 ---
% 官方标准红 RGB: 193, 0, 22 (大致)
% 为了简历阅读舒适，这里稍微调暗一点点，显得稳重
\definecolor{HUST_Red}{RGB}{176, 31, 36} 
% 如果你喜欢原来那种蓝色风格，也可以保留蓝色，华科理工科用深蓝也很合适
% \definecolor{HUST_Blue}{RGB}{0, 51, 102} 

% --- 标题格式设置 (使用 HUST_Red) ---
\titleformat{\section}
  {\LARGE\bfseries\raggedright}
  {}{0em}
  {}
  [{\color{HUST_Red}\titlerule}] % 标题下划线颜色
\titlespacing*{\section}{0cm}{*1.2}{*1.2}

\titleformat{\subsection}
  {\large\bfseries\raggedright}
  {}{0em}
  {}
  []
\titlespacing*{\subsection}{0cm}{*1.2}{*1.2}

% 页面大小与页边距
\usepackage[
	a4paper,
	left=1.2cm,
	right=1.2cm,
	top=1.8cm, % 稍微增加顶部距离以容纳页眉
	bottom=1cm,
	nohead
]{geometry}

\renewcommand{\arraystretch}{1.2}
\linespread{1.2}
\renewcommand{\CJKglue}{\hskip 0.05em}

% --- 字体设置 (优化版) ---
% 原模板使用了特定路径的字体，如果在 Overleaf 或没有该字体会报错。
% 这里做了判断：如果能找到 SweiSpring 就用，找不到就用系统自带的。

\ifxetex
  % 尝试加载自定义字体，如果不想用自定义字体，直接注释掉下面这块
  % \setmainfont[Path=fonts/,Extension=.ttf,BoldFont=* Bold]{SweiSpring}
  % \setCJKmainfont[Path=fonts/,Extension=.ttf,BoldFont=* Bold]{SweiSpring}
  
  % --- 推荐方案：使用通用字体 (Overleaf 上可用 Fandol 系列) ---
  % 如果你在本地 Windows，可以改为 SimSun (宋体) 和 SimHei (黑体)
  \setCJKmainfont[BoldFont=FandolHei-Regular, ItalicFont=FandolKai-Regular]{FandolSong-Regular}
  \setmainfont{Times New Roman} % 英文用 Times 比较正式
\fi

% 填写个人信息
% 学院 / 部门
\newcommand{\school}{人工智能与自动化学院 | School of AIA} 

% 联系方式
\newcommand{\contact}{
    \scriptsize
    \textcolor{white}{
        % 邮箱 (改为 hust 后缀)
        \faEnvelope \quad \href{mailto:youremail@hust.edu.cn}{youremail@hust.edu.cn}
        \hspace{4em}
        % 手机号
        \faPhone \quad  130-6666-8888
        \hspace{4em}
        % 微信/Github
        \faGithub \quad \href{https://github.com/yourname}{github.com/yourname}
    }
}

%%%%%%%%%%%%%%%%%%%%
% 简历正文
%%%%%%%%%%%%%%%%%%%%
\begin{document}

    % ============================================================
    % 页眉设计 (自动绘制红色背景条，无需 header.png)
    % ============================================================
    \begin{tikzpicture}[remember picture, overlay]
        % 1. 画红色背景条
        \fill[HUST_Red] (current page.north west) rectangle ([yshift=-1.6cm]current page.north east);
        
        % 2. 放置左侧校徽文字 Logo (请确保 images/hust_logo.png 是白字或透明底)
        % 如果没有白字Logo，可以将下面这行注释掉，直接用文字写校名
        \node[anchor=west] at ([xshift=1cm, yshift=-0.8cm]current page.north west){
            \includegraphics[height=1.0cm]{images/hust_logo.pdf} 
        };
        
        % (备用方案：如果没有图片Logo，用文字代替，取消下面注释)
        % \node[anchor=west] at ([xshift=1cm, yshift=-0.8cm]current page.north west){
        %     \color{white}\Huge\bfseries 华中科技大学
        % };

        % 3. 放置右侧学院名称
        \node[anchor=east] at ([xshift=-1cm, yshift=-0.8cm]current page.north east){
            \textcolor{white}{\textbf{\school}}
        };
    \end{tikzpicture}
    \vspace{-2em} % 调整正文与页眉的距离

    % ============================================================
    % 页脚设计 (自动绘制红色背景条，无需 footer.png)
    % ============================================================
    \begin{tikzpicture}[remember picture, overlay]
        % 画底部红色背景条
        \fill[HUST_Red] (current page.south west) rectangle ([yshift=1.2cm]current page.south east);
        
        % 放置联系方式
        \node[anchor=center] at ([yshift=0.6cm]current page.south){
            \contact
        };
    \end{tikzpicture}

    % ============================================================
    % 背景水印 (华科校徽)
    % ============================================================
    \begin{tikzpicture}[remember picture, overlay]
        \node[
            opacity=0.1,           % 透明度：建议 0.1 到 0.15 之间 (太深会干扰文字)
            blend mode=multiply,    % 混合模式：让图片的白色背景变透明 (类似PS的正片叠底)
            inner sep=0pt           % 去除多余边距
        ] at(current page.center){
            \includegraphics[width=0.6\paperwidth, keepaspectratio]{images/hust_badge.png}
        };
    \end{tikzpicture}

    % ============================================================
    % 个人信息板块
    % ============================================================
    \begin{figure}[h]
        % 左半边：文字信息
        \begin{minipage}{0.82\textwidth}
            \section{\makebox[\widthof{\faAddressCard}][c]{\color{HUST_Red}{\faAddressCard}}\quad 个人信息}
            \begin{tabularx}{\linewidth}{p{\widthof{出生日期:}}Xp{\widthof{政治面貌:}}X}
                姓\ \ \ \ \ \ \ \ 名: & 喻家山居士 & 
                性\ \ \ \ \ \ \ \ 别: & 男  \\
                出生年月: & 2000年1月 & 
                政治面貌: & 群众 \\
                意向岗位: & 算法工程师 &
                籍\ \ \ \ \ \ \ \ 贯: & 湖北武汉
            \end{tabularx}
        \end{minipage}
        \hspace{1em}
        % 右半边：照片
        \begin{minipage}{0.14\textwidth}
            \setlength{\fboxsep}{0pt}
            % 边框颜色改为华科红
            \fcolorbox{HUST_Red}{white}{\includegraphics[width=\linewidth]{images/avatar.png}}
        \end{minipage}
    \end{figure}
    \vspace{-1.5em}

    % ============================================================
    % 教育背景
    % ============================================================
    \section{\makebox[\widthof{\faGraduationCap}][c]{\color{HUST_Red}{\faGraduationCap}}\quad 教育背景}

    % 示例内容
    {\large \textbf{华中科技大学}}，本科 \hfill {湖北，武汉} \\
    {{启明学院}}，专业：自动化 \hfill {2018年9月-2022年6月} \\
    {主修课程}：自动控制原理、信号与系统、数字图像处理、C++程序设计\ 等。

    \vspace{0.5em}
    {\large \textbf{华中科技大学}}，硕士 \hfill {湖北，武汉} \\
    {{人工智能与自动化学院}}，专业：控制科学与工程 \hfill {2022年9月-至今} \\
    \textbf{研究方向}：计算机视觉、多模态大模型\ 等。

    % ============================================================
    % 科研成果
    % ============================================================
    \section{\makebox[\widthof{\faBook}][c]{\color{HUST_Red}{\faBook}}\quad 科研成果}

    \textbf{Paper Title Example 1.} \\
    \textbf{San Zhang (Your Name)}, Si Li, Wu Wang. \hfill 
    发表于 \textbf{CVPR 2024} (CCF-A类会议)

    \vspace{0.5em}
    \textbf{Paper Title Example 2.} \\
    Si Li, \textbf{San Zhang}, Wu Wang. \hfill 
    发表于 \textbf{IEEE TPAMI} (SCI 一区, IF=20.1)

    % ============================================================
    % 项目与实习
    % ============================================================
    \section{\makebox[\widthof{\faLaptopCode}][c]{\color{HUST_Red}{\faLaptopCode}}\quad 项目与实习经历}
    
    \vspace{0.5em}
    {\large{\textbf{某某大厂核心研发部}}}，算法实习生 \hfill {2023年6月 - 2023年9月}\\
    \textbf{主要职责}：负责多模态图文检索模型的预训练与微调。\\
    \textbf{工作产出}：优化了特征提取模块，使下游任务准确率提升了 2.5\%。


    % ============================================================
    % 技能与证书
    % ============================================================
    \section{\makebox[\widthof{\faWrench}][c]{\color{HUST_Red}{\faWrench}}\quad 技能与证书}
    \vspace{0.5em}
    \begin{itemize}
        \item \textbf{编程语言}：熟练掌握 Python, C++, MATLAB; 熟悉 Linux 开发环境。
        \item \textbf{开发框架}：精通 PyTorch, TensorFlow, TensorRT 推理加速。
        \item \textbf{语言能力}：CET-6 (600+), 托福 (105), 能流畅阅读英文文献。
        \item \textbf{荣誉奖项}：国家奖学金 (2021), "互联网+"大赛国家银奖 (2022).
    \end{itemize}

\end{document}